\documentclass[12pt]{article}

\usepackage{amsmath, amssymb, amsthm}
\usepackage{geometry}
\geometry{margin=1in}

\title{LECTURE SCRIBE}
\author{CSE400 \\ Dhaval Patel, PhD}
\date{February 10, 2026}

\begin{document}

\maketitle

\section*{Transformation of RVs + Function of Two RVs + Example}

\section{Transformation of Random Variables}

\subsection{Lecture Outline (as presented on slide)}

\begin{enumerate}
    \item \textbf{Transformation of Random Variables} \\
    Learning of transformation techniques for random variables.
    
    \item \textbf{Function of Two Random Variables} \\
    Joint transformations and derived distributions.
    
    \item \textbf{Illustrative Example} \\
    Detailed derivation for the case:
    \[
    Z = X + Y
    \]
\end{enumerate}

\subsection{Key Objective}

\[
\boxed{\text{How to find PDF of new transformed RV, } f_Y(y)\text{?}}
\]

\subsection{Assumption}

\[
\textbf{Assumption: } f_X(x) \rightarrow \text{PDF of } X \text{ is known a priori}
\]

Thus, throughout the transformation procedure, \( f_X(x) \) is assumed known.

\section{Function of One Random Variable}

\subsection{Definition}

Let:
\[
Y = g(X)
\]

The slide explicitly separates two cases:

\begin{itemize}
    \item (+ve) Monotonically increasing
    \item (–ve) Monotonically decreasing
\end{itemize}

Both cases are handled separately.

\section{Case 1: \( g(x) \) Monotonically Increasing}

\subsection{Step S1 – Compute CDF}

\[
F_Y(y) = \Pr(Y \le y)
\]

Since \( Y = g(X) \):

\[
= \Pr(g(X) \le y)
\]

Because \( g(x) \) is increasing:

\[
= \Pr(X \le g^{-1}(y))
\]

Therefore:

\[
F_Y(y) = F_X(g^{-1}(y))
\]

\subsection{Step S2 – Differentiate w.r.t. \( y \)}

\[
f_Y(y) = \frac{d}{dy}\left[F_X(g^{-1}(y))\right]
\]

\[
= f_X(g^{-1}(y)) \cdot \frac{d}{dy}g^{-1}(y)
\]

The slide also writes the equivalent form:

\[
f_Y(y) = f_X(x)\left|\frac{dx}{dy}\right|_{x=g^{-1}(y)}
\]

\section{Case 2: \( g(x) \) Monotonically Decreasing}

\subsection{Step S1 – Compute CDF}

\[
F_Y(y) = \Pr(Y \le y)
\]

\[
= \Pr(g(X) \le y)
\]

Since \( g(x) \) is decreasing:

\[
= \Pr(X \ge g^{-1}(y))
\]

\[
= 1 - F_X(g^{-1}(y))
\]

\subsection{Step S2 – Differentiate}

\[
\boxed{
f_Y(y)
=
\frac{f_X(x)}{\left|\frac{dy}{dx}\right|}_{x=g^{-1}(y)}
}
\]

\subsection{Step S3}

Change the limits for \( y \).

\section{Example: Single Variable Transformation}

\subsection{Given}

\[
X \sim \text{Uniformly distributed RV } \sim (-1,1)
\]

\[
Y = g(X) = \sin\left(\frac{\pi X}{2}\right)
\]

The PDF of \( X \) is given:

\[
f_X(x)=
\begin{cases}
\frac{1}{2}, & -1 < x < 1 \\
0, & \text{otherwise}
\end{cases}
\]

\subsection{Solution}

From:

\[
y = \sin\left(\frac{\pi x}{2}\right)
\]

Invert:

\[
x = \frac{2}{\pi}\sin^{-1}(y)
\]

Differentiate:

\[
\frac{dx}{dy}
=
\frac{2}{\pi}
\cdot
\frac{1}{\sqrt{1-y^2}}
\]

Using transformation formula:

\[
f_Y(y)=f_X(x)\left|\frac{dx}{dy}\right|
\]

Substitute:

\[
\frac{1}{2}
\cdot
\frac{2}{\pi}
\cdot
\frac{1}{\sqrt{1-y^2}}
\]

\[
\frac{1}{\pi\sqrt{1-y^2}}
\]

Support shown:

\[
-1<y<1
\]

\section{Function of Two Random Variables}

\subsection{Setup}

\[
Z = X + Y
\]

The slide lists:

\begin{enumerate}
    \item Find PDF of \( Z \), \( f_Z(z) \)
    \item \( f_Z(z) \), if \( X \) \& \( Y \) are independent
    \item Let \( X \sim N(0,1) \) \& \( Y \sim N(0,1) \), Prove \( Z \sim N(0,2) \)
    \item If \( X \) \& \( Y \) are exponential distribution RVs with parameter \( \lambda \), find \( f_Z(z) \)
\end{enumerate}

\section{Derivation of PDF of \( Z = X + Y \)}

\subsection{Step 1 – CDF}

\[
F_Z(z)=\Pr(Z \le z)
\]

\[
=\Pr(X+Y \le z)
\]

Region from slide:

\[
x \le z-y
\]

Thus:

\[
F_Z(z)
=
\int_{-\infty}^{\infty}
\int_{-\infty}^{z-y}
f_{XY}(x,y)\,dx\,dy
\]

\section{Leibniz Rule (As Presented)}

If:

\[
G(x)=\int_{a(x)}^{b(x)} h(x,y)\,dy
\]

Then:

\[
\frac{d}{dx}G(x)
=
\frac{d}{dx}b(x)\,h(x,b(x))
-
\frac{d}{dx}a(x)\,h(x,a(x))
+
\int_{a(x)}^{b(x)}
\frac{\partial}{\partial x}h(x,y)\,dy
\]

The slide identifies:

\begin{itemize}
    \item Upper limit term
    \item Lower limit term
    \item Integral term
\end{itemize}

\section{PDF of \( Z \)}

\[
f_Z(z)=\frac{d}{dz}F_Z(z)
\]

Applying the rule as shown:

\[
f_Z(z)
=
\int_{-\infty}^{\infty}
f_{XY}(z-y,y)\,dy
\]

Equivalent vertical strip form:

\[
\boxed{
f_Z(z)
=
\int_{-\infty}^{\infty}
f_{XY}(x,z-x)\,dx
}
\]

\section{Independent Case}

If \( X \) \& \( Y \) are independent:

\[
f_{XY}(x,y)=f_X(x)f_Y(y)
\]

Thus:

\[
\boxed{
f_Z(z)
=
\int_{-\infty}^{\infty}
f_X(x)f_Y(z-x)\,dx
}
\]

Convolution Integral

\section{Normal Example}

Given:

\[
X \sim N(0,1), \quad Y \sim N(0,1)
\]

\[
f_X(x)=\frac{1}{\sqrt{2\pi}}e^{-x^2/2}
\]

\[
f_Y(y)=\frac{1}{\sqrt{2\pi}}e^{-y^2/2}
\]

Using convolution:

\[
f_Z(z)
=
\int_{-\infty}^{\infty}
\frac{1}{\sqrt{2\pi}}e^{-x^2/2}
\frac{1}{\sqrt{2\pi}}e^{-(z-x)^2/2}
dx
\]

Algebra shown leads to:

\[
\frac{1}{\sqrt{2\pi\cdot2}}
e^{-z^2/(2\cdot2)}
\]

Therefore:

\[
\boxed{Z \sim N(0,2)}
\]

\section{Exponential Example}

Given:

\[
f_X(x)=\lambda e^{-\lambda x}, \quad x>0
\]

\[
f_Y(y)=\lambda e^{-\lambda y}, \quad y>0
\]

Using convolution:

\[
f_Z(z)
=
\int_{-\infty}^{\infty}
\lambda e^{-\lambda x}
\lambda e^{-\lambda(z-x)}
u(x)u(z-x)\,dx
\]

Slide specifies:

\[
u(x)u(z-x)
=
\begin{cases}
1, & 0<x<z \\
0, & \text{otherwise}
\end{cases}
\]

Thus:

\[
f_Z(z)
=
\lambda^2 e^{-\lambda z}
\int_0^z dx
\]

\[
\lambda^2 z e^{-\lambda z}
\]

\section{Example: \( Z = X - Y \)}

\[
F_Z(z)=\Pr(Z \le z)
\]

\[
=\Pr(X-Y \le z)
\]

Region shown:

\[
x \le z+y
\]

Thus:

\[
F_Z(z)
=
\int_{-\infty}^{\infty}
\int_{-\infty}^{z+y}
f_{XY}(x,y)\,dx\,dy
\]

\[
f_Z(z)=\frac{d}{dz}F_Z(z)
\]

Same procedure, \( Z = X + Y \).

\section{Example: Ratio}

\[
Z=\frac{X}{Y}
\]

\[
F_Z(z)=\Pr\left(\frac{X}{Y} \le z\right)
\]

The slide distinguishes:

\begin{itemize}
    \item \( Y>0 \)
    \item \( Y<0 \)
\end{itemize}

Region drawn accordingly.

\section{Example: Radial Transformation}

\[
R=\sqrt{X^2+Y^2}
\]

\[
f_R(r)=?
\]

\end{document}
